\chapter{Introduction}
\label{cha:introduction}

\section{Motivation and Problem statement}
\label{sec:intro-problem}

Policies for deformable manipulation are learned from recordings of real interactions, and
digitising those recordings by hand takes 3D modelling, material parameters and manual matching. The goal is to automate that process, converting multiple RGB video recordings of the same interaction into a single \emph{digital twin} of the manipulated object.

\textcite{phystwin} proposed PhysTwin, a real2sim pipeline that takes RGBD video from several
cameras, segments the object, tracks the object and the controller, and fits a spring--mass
\emph{warp} model for interaction in NVIDIA Warp \parencite{warp}. This work improves the robustness of PhysTwin by (1) removing the depth sensor, (2) removing camera calibration and (3) tolerating imperfect segmentation and tracking; the resulting pipeline is called DIPhysTwin. It should allow PhysTwin to be used in more flexible capture conditions, with varying degrees of camera setup and less human intervention.

\section{Contributions}
\label{sec:intro-contributions}

The main contributions of this work are:

\begin{enumerate}
	\item \textbf{A depth-inference-first pipeline.} Poses and metric scale come from a dedicated
	stage rather than the depth model, and depth is estimated per view.
	\item \textbf{A ground-truth-free temporal depth smoother.} A small network learns a
	\emph{residual correction} to one depth model's maps, smoothing temporal flicker in depth videos while preserving spatial consistency.
	\item \textbf{Robustness engineering.} The pipeline handles resolutions independently, runs
	idempotently, aligns the scene to the ground with RANSAC and propagates that transform to poses,
	tracks and meshes, with fallbacks when shape prior fitting or dense tracking fails.
\end{enumerate}


